\chapter{Decode Stage}

\section{Overview}

The Decode stage receives fetched instructions from Fetch-2 and converts them into an internal
representation suitable for later pipeline stages.

Its responsibilities include:
\begin{itemize}
  \item decoding the instruction opcode,
  \item extracting source and destination registers,
  \item generating immediate values,
  \item identifying instruction formats,
  \item producing control information,
  \item forwarding decoded instructions to the Rename stage.
\end{itemize}

This stage forms the interface between instruction fetch and out-of-order execution.

\begin{rtlrefs}
  \rtlitem{rtl/datapath/rtl/id_stage/rtl/decoder.sv}
  \rtlitem{rtl/datapath/rtl/id_stage/rtl/immediate.sv}
\end{rtlrefs}

\begin{figure}[htbp]
  \centering
  \includegraphics[width=0.28\textwidth]{ch06_decode_overview.png}
  \caption{Decode stage: F2-to-D $\rightarrow$ Decoder $\rightarrow$ Instruction Queue}
  \label{fig:decode-overview}
\end{figure}

\begin{figure}[htbp]
  \centering
  \includegraphics[width=0.26\textwidth]{ch06_decode_detail.png}
  \caption{Detailed signal view of the Decode stage}
  \label{fig:decode-detail}
\end{figure}

\section{Instruction Decode}

The Decoder interprets the binary instruction according to the RISC-V ISA.
It identifies:
\begin{itemize}
  \item instruction format,
  \item functional unit,
  \item register operands,
  \item immediate usage,
  \item memory operations,
  \item branch instructions,
  \item CSR accesses,
  \item floating-point instructions,
  \item vector instructions.
\end{itemize}

The decoded information is forwarded toward Rename.

\begin{rtlrefs}
  \rtlitem{rtl/datapath/rtl/id_stage/rtl/decoder.sv}
\end{rtlrefs}

\section{Immediate Generation}

Many RISC-V instructions contain embedded immediate fields.
The Immediate Generator extracts these fields and reconstructs complete signed values
according to the instruction format.

Supported formats include:
\begin{itemize}
  \item I-type
  \item S-type
  \item B-type
  \item U-type
  \item J-type
\end{itemize}

These immediates are later consumed by the ALU, Branch Unit, and Address Generation Unit.

\begin{rtlrefs}
  \rtlitem{rtl/datapath/rtl/id_stage/rtl/immediate.sv}
\end{rtlrefs}

\section{Opcode Decoding}

Opcode decoding determines the exact class of instruction.
Examples include:
\begin{itemize}
  \item Arithmetic
  \item Logical
  \item Branch
  \item Jump
  \item Load
  \item Store
  \item Floating Point
  \item SIMD
  \item CSR
  \item System
\end{itemize}

The Decoder produces internal control signals used by downstream stages to select the
appropriate execution unit.

\begin{rtlrefs}
  \rtlitem{rtl/datapath/rtl/id_stage/rtl/decoder.sv}
\end{rtlrefs}

\section{Control Signal Generation}

Based on the decoded instruction, the Decode stage generates the control information required
throughout the remainder of the pipeline.

These signals specify:
\begin{itemize}
  \item destination register validity,
  \item operand selection,
  \item ALU operation,
  \item branch type,
  \item memory access type,
  \item CSR operation,
  \item vector operation,
  \item floating-point operation.
\end{itemize}

These control signals accompany the instruction into the Rename stage.

\begin{rtlrefs}
  \rtlitem{rtl/datapath/rtl/id_stage/rtl/decoder.sv}
  \rtlitem{rtl/control_unit/rtl/control_unit.sv}
\end{rtlrefs}

\section{Output to Rename Stage}

After decoding completes, the instruction is forwarded to the Rename stage.

At this point the instruction has:
\begin{itemize}
  \item decoded opcode,
  \item source registers,
  \item destination register,
  \item immediate value,
  \item generated control signals.
\end{itemize}

The Rename stage will subsequently eliminate false register dependencies using register
renaming before dispatching instructions toward execution.

\begin{rtlrefs}
  \rtlitem{rtl/datapath/rtl/id_stage/rtl/decoder.sv}
  \rtlitem{rtl/datapath/rtl/ir_stage/rtl/rename_table.sv}
  \rtlitem{rtl/datapath/rtl/ir_stage/rtl/instruction_queue.sv}
\end{rtlrefs}

\begin{figure}[htbp]
  \centering
  \includegraphics[width=0.55\textwidth]{ch06_decode_to_rename.png}
  \caption{Decoded instructions flowing from the Decode stage into the Rename stage}
  \label{fig:decode-to-rename}
\end{figure}

\subsection*{Example (adapted from \texttt{doc/specification/decoder\_specification})}

The cases are:

\begin{itemize}
  \item Cycle 1: Invalid input $\rightarrow$ invalid output.
  \item Cycle 2: Exception input $\rightarrow$ exception output and not decoded.
  \item Cycle 3: Regular case, valid input, not illegal instruction, and signed operation.
  \item Cycle 4: Regular case again, but it has the power to change the PC.
  \item Cycle 5: Case with illegal instruction, hence exception valid to 1.
  \item Cycle 6: Case of the fence with the special signal to 1.
\end{itemize}

In the following figure we can see all the cases mentioned above.

\begin{figure}[htbp]
  \centering
  \includegraphics[width=\textwidth]{ch06_decoder_waveform.png}
  \caption{Decoder timing waveform illustrating the six example cases}
  \label{fig:decoder-waveform}
\end{figure}
